\section{Building and verifying the formalization}\label{app:verification}

The Lean project is the directory \leanfile{lean/} next to this paper. It uses Lean and Mathlib, both at version v4.34.1
(\leanfile{lean-toolchain}, \leanfile{lakefile.toml}, \leanfile{lake-manifest.json}), and requires the shared
\leanfile{lean-machine/} library of the repository through a relative path, which supplies Mathlib at the pinned
revision. It uses no library of an earlier paper. Keep the repository's directory layout so that
\leanfile{lean-machine/} stays in place.

\paragraph{Build.} With \texttt{elan} installed, run in \leanfile{lean/}:
\begin{center}
\texttt{lake exe cache get \&\& lake build}
\end{center}
The first command downloads the compiled Mathlib cache; the second compiles \leanfile{SwapRigidity}.

\paragraph{Axioms and index.} The file \leanfile{lean/PaperAxioms.lean}, generated from the \LaTeX{} sources of this paper
by \leanfile{tools/lean_index.py}, runs \texttt{\#print axioms} on every declaration cited in this paper
(\cref{app:leanindex}):
\begin{center}
\texttt{lake env lean PaperAxioms.lean}
\end{center}
Every declaration depends at most on \texttt{propext}, \texttt{Classical.choice} and \texttt{Quot.sound}. No Lean file
of this library uses \texttt{sorry}, \texttt{admit} or \texttt{native\_decide}, declares an \texttt{axiom}, bypasses the
kernel via \texttt{unsafe}, \texttt{implemented\_by} or \texttt{extern}, or raises \texttt{maxHeartbeats}. Running
\texttt{python3 tools/lean\_index.py} from the paper's root directory regenerates both
\leanfile{paper/sections/leanindex.tex} and \leanfile{lean/PaperAxioms.lean} (\texttt{--index-only} writes only the
index; \texttt{--check} verifies that every cited name resolves and that both generated files are current).

\paragraph{Where things are.} The library \leanfile{lean/SwapRigidity/} has the following modules, all imported by the
root module \leanfile{lean/SwapRigidity.lean}.
\begin{itemize}[leftmargin=1.5em]\raggedright
\item \leanfile{Flag.lean} (\cref{sec:order}): the swap code, the elimination on staircase patterns, invariance under
changes of bases, the coordinate form of Theorem~FL, the corollary for a pair of flags, the exchange of sides, and
attainment.
\item \leanfile{Pigeonhole.lean} (\cref{sec:symmetry}): the pigeonhole lemma in its three forms, the rank bound for
commutators and its subadditivity under composition, and the dimension count for an intersection of preimages.
\item \leanfile{Swap.lean}: the dimensions of the kernel and image of $\sw-1$ and the bound $\dim\ker[\sw,T]\ge d$.
\item \leanfile{ClassUniform.lean}: class-uniform codes and Theorem~TP on a cycle.
\item \leanfile{Entry.lean} (\cref{prop:entry}): rules given entry by entry, and the bound for decoders that change at
single entries.
\item \leanfile{EntryAdvice.lean} (\cref{cor:entry-advice}): the same rules with advice of any kind, over a finite
field.
\item \leanfile{BlockSymmetric.lean}: the block-symmetric code, its dimension, and attainment of the bound.
\item \leanfile{Line.lean}: two-sided rules on a line and the line form of Theorem~TP.
\item \leanfile{Twist.lean}: the twist-class theorem.
\item \leanfile{AperiodicExample.lean} and the generated \leanfile{AperiodicExampleData.lean} (\cref{prop:aperiodic}):
an explicit class-uniform code of dimension $25$ for $d=6$ over $\F_2$. The data file contains definitions only. The
finite facts about the data are proved by \texttt{decide}, which the kernel checks; nothing is taken on trust from the
certificate \leanfile{tools/data/aperiodic_d6_dim25.json} or from the generator
\leanfile{tools/gen_aperiodic_example.py} (run it with \texttt{--check} to confirm that the data file is current).

\end{itemize}

\paragraph{Computations.} Every script that was used for this paper is in \leanfile{tools/}, and the saved codes are
in \leanfile{tools/data/}. The file \leanfile{tools/README.md} lists the scripts with their commands and the versions
used, and records the solver runs behind \cref{tab:small}, including those that timed out. Each script documents its
usage at its top.
\begin{itemize}[leftmargin=1.5em]\raggedright
\item \leanfile{tools/exhaustive_flag_checks.py} and \leanfile{tools/exhaustive_pigeonhole_checks.py} (Python with
NumPy): the exhaustive and randomized checks behind the numerical remarks of \cref{sec:order,sec:symmetry} and the
entries for $d\le4$ of \cref{tab:small}.
\item \leanfile{tools/three_term_example.py}: the examples of \cref{rem:three}.
\item \leanfile{tools/digit_and_composition_checks.py}: small instances of \cref{cor:digit} and of
\cref{rem:compose-TP}, both of which are proved in the text.
\item \leanfile{tools/entry_form_check.py}: small instances of \cref{prop:entry}, which is proved in Lean apart from
a count.
\item \leanfile{tools/two_class_family.py}: the rule of \cref{prop:two-class}, which is proved in the text, over five
small fields, with every rank stated in the proof; the saved codes over $\F_2$; and small instances of
\cref{cor:every-M} (\cref{rem:two-class}).
\item \leanfile{tools/class_uniform_sat.py} (needs the package \texttt{pycryptosat}): one run of the solver, for a
partition and a target dimension.
\item \leanfile{tools/small_sizes.py}: \cref{tab:small} with its bounds, recomputed, and with every saved code checked
again; the list of partitions up to affine maps; the search that raises the target dimension, which produced the saved
codes of the table; and controls that repeat the answers of the solver quoted in \cref{rem:small}.
\item \leanfile{tools/verify_class_uniform_code.py}: checks every saved code in \leanfile{tools/data/} against the
definition.
\item \leanfile{tools/analyse_code.py}: the intersections of the code of \cref{prop:aperiodic} with its translates,
quoted after that proposition.
\end{itemize}
Apart from the certificate of \cref{prop:aperiodic}, which Lean checks, these computations are not part of the formal
development, and no proof depends on them.
